\chapter{Model Training and MLflow}
\index{MLflow}\index{experiments}\index{model registry}\index{mlflow.autolog}\index{hyperparameter tuning}\index{model versioning}\index{Random Forest}\index{Scikit-Learn}%

\begin{objectives}
\begin{itemize}
    \item Explain Fabric's machine-learning items: experiments, models, and their MLflow integration.
    \item Track parameters, metrics, and artifacts with \texttt{mlflow.autolog()} and explicit logging.
    \item Register and version models in the Fabric model registry.
    \item Train and auto-track a Scikit-Learn Random Forest, registering the best run.
    \item Compare hundreds of hyperparameter runs efficiently in a shared workspace.
    \item Attach the metadata---dataset version, code commit---that makes a model auditable and rollback-able.
\end{itemize}
\end{objectives}

\begin{crosscloud}
Experiment tracking and model registries are the two pieces of ML infrastructure every platform converged on, and MLflow is the closest thing the industry has to a common denominator---it is native in Fabric and Databricks and deployable everywhere else.

\vspace{2mm}
{\footnotesize
\begin{tabular}{@{}p{2.8cm}p{3.2cm}p{3.2cm}p{3.2cm}p{3.2cm}@{}}
\toprule
\textbf{Capability} & \textbf{Fabric} & \textbf{AWS} & \textbf{GCP} & \textbf{Snowflake} \\
\midrule
Experiment tracking & Built-in MLflow & SageMaker Experiments / MLflow on EC2 & Vertex AI Experiments & MLflow (external or integrated) \\
Model registry & Fabric model registry (MLflow) & SageMaker Model Registry & Vertex AI Model Registry & Snowflake Model Registry \\
Autologging & \texttt{mlflow.autolog()} & Same library & Same library & Same library \\
Training compute & Fabric Spark pools & SageMaker Training / EMR & Vertex Training / Dataproc & Snowpark ML / warehouses \\
Artifact store & OneLake (via MLflow) & S3 & Cloud Storage & Stages / tables \\
Run comparison UI & Experiment item UI & SageMaker Studio & Vertex console & Registry UI / notebooks \\
\bottomrule
\end{tabular}}

\vspace{1mm}
Databricks' managed MLflow is the direct ancestor of Fabric's integration. MongoDB does not host training; it supplies features. The portable claim of this chapter: \emph{if it is not tracked, it did not happen; if it is not registered, it cannot be governed; if it cannot be rolled back, it is not production.}
\end{crosscloud}

\begin{keyterms}
\textbf{Experiment}: the MLflow container grouping related training runs. \quad
\textbf{Run}: one execution of training code with its parameters, metrics, and artifacts. \quad
\textbf{\texttt{mlflow.autolog()}}: automatic capture of framework-specific parameters, metrics, and models per run. \quad
\textbf{Model registry}: the governed catalog of versioned, stage-tagged models. \quad
\textbf{Model version}: an immutable, numbered registration of a trained model. \quad
\textbf{Alias/stage}: a movable pointer (for example \texttt{Production}) over versions enabling rollback.
\end{keyterms}

\section{Week 18 Roadmap}
Week~18 makes machine learning auditable. Monday covers experiments and models as Fabric items. Tuesday is MLflow tracking: parameters, metrics, artifacts, and autolog. Wednesday covers the model registry: registration, versioning, and aliases. Thursday is the hands-on project: train a Random Forest with \texttt{mlflow.autolog()}, track it, register the best run. Friday designs the collaborative workflow for hundreds of hyperparameter runs in a shared workspace.

Chapter~17 made features honest; this week makes models honest. The untracked notebook that produced the production model is organizational debt with interest: nobody can reproduce it, improve it, or roll back from it \cite{microsoftMlflowFabric,mlflowDocs}.

\begin{notebox}
The minimum metadata contract for every training run: parameters, metrics, dataset version (or as-of date), code commit, and the environment. Autolog gives you the first three cheaply; the last two are discipline.
\end{notebox}

\section{Monday: Experiments and Models in Fabric}
\index{experiments!Fabric items}\index{models!Fabric items}
Fabric surfaces MLflow through two workspace items \cite{microsoftMlflowFabric}. An \textbf{experiment} item is the tracking container: open it to browse runs, compare metrics, and inspect artifacts. A \textbf{model} item is the registry entry: versioned, aliased, and loadable by name from any notebook.

\begin{definitionbox}
An \textbf{MLflow experiment} groups related training runs under one name; each \textbf{run} records the parameters, metrics, tags, and artifacts of one execution, so work is comparable, searchable, and reproducible.
\end{definitionbox}

The workspace integration changes the collaboration physics: experiments are visible to the whole workspace (not trapped on one laptop), runs accumulate as organizational memory, and the registry attaches governance---who may promote a model to \texttt{Production} is a permission, not a convention.

\section{Tuesday: Tracking with MLflow}
\index{MLflow!tracking}\index{mlflow.autolog}
\subsection{Autolog}
One line before training turns a Scikit-Learn, Spark MLlib, or XGBoost fit into a fully tracked run: parameters, metrics, and the model artifact itself \cite{mlflowDocs,microsoftMlflowFabric}.

\begin{lstlisting}[language=Python,caption={Autologged training run with explicit metadata tags.},label={lst:chapter18-autolog}]
import mlflow
from sklearn.ensemble import RandomForestClassifier
from sklearn.model_selection import train_test_split
from sklearn.metrics import roc_auc_score

mlflow.set_experiment("churn-model-week18")
mlflow.autolog()

df = spark.read.table("gold.churn_features").where("as_of_date = '2026-07-31'").toPandas()
X = df.drop(columns=["customer_id", "churned", "as_of_date"])
y = df["churned"]
X_train, X_test, y_train, y_test = train_test_split(
    X, y, test_size=0.2, random_state=42, stratify=y)

with mlflow.start_run():
    mlflow.set_tag("dataset_as_of", "2026-07-31")
    mlflow.set_tag("code_commit", "a1b2c3d")        # from your CI or notebook env
    model = RandomForestClassifier(n_estimators=200, max_depth=12, random_state=42)
    model.fit(X_train, y_train)
    auc = roc_auc_score(y_test, model.predict_proba(X_test)[:, 1])
    mlflow.log_metric("holdout_auc", auc)
\end{lstlisting}

\subsection{Explicit Logging}
Autolog captures the framework's view; you log the rest: business metrics (expected revenue saved per flagged churner), data slices (AUC by region), and environment facts. A run without its business metric is a science project; with it, a decision input.

\begin{tipbox}
Log one ``headline'' metric with a stable name across every experiment in the program. Cross-experiment comparison---this month versus last, this model family versus that---requires a metric that means the same thing everywhere.
\end{tipbox}

\section{Wednesday: The Model Registry}
\index{model registry}\index{model versioning}\index{aliases}
Registration promotes a run's artifact into the governed catalog: the model gets a name, an immutable version number, and movable aliases such as \texttt{Staging} and \texttt{Production} \cite{microsoftMlflowFabric,mlflowDocs}.

\begin{definitionbox}
The \textbf{model registry} is the versioned catalog of trained models; an \textbf{alias} is a movable pointer over versions, so promoting or rolling back production is a pointer move, not a redeploy.
\end{definitionbox}

Why register rather than reference run artifacts directly: the registry is the governance boundary. Versions are immutable (the artifact a version pointed at never changes), aliases decouple consumers from versions (scoring code loads \texttt{models:/churn\_model/Production} and never hardcodes a version), and the audit trail---who promoted what, when, with which metrics---lives in one place.

\begin{pitfall}
Scoring code that loads a specific run ID or version number bypasses the entire governance point. Consumers load by alias; humans move aliases under review; rollbacks are alias moves. The day a hardcoded version ships to production, the registry becomes decorative.
\end{pitfall}

\section{Thursday: Hands-on Project}
\index{hands-on lab}\index{MLflow!hands-on}
The Week~18 lab trains a Scikit-Learn Random Forest on the Chapter~17 feature table, tracks the run with \texttt{mlflow.autolog()}, and registers the best model.

\subsection{Goal and Architecture}
Reuse \texttt{gold.churn\_features} (or a provided classification dataset) in workspace \texttt{fabric-de-week18}. Run a small hyperparameter grid (six runs), compare runs in the experiment UI, and register the winner as version~1 of \texttt{churn\_model}.

\subsection{Step-by-Step Actions}
\begin{enumerate}
    \item Set the experiment and enable autolog.
    \item Load one as-of partition of the feature table; split with stratification and a fixed seed.
    \item Train six configurations across \texttt{n\_estimators} and \texttt{max\_depth}, each in its own run with the metadata tags.
    \item Log \texttt{holdout\_auc} explicitly in every run.
    \item In the experiment item, sort by \texttt{holdout\_auc}; inspect the best run's parameters and artifacts.
    \item Register the best run's model as \texttt{churn\_model} version~1; assign the \texttt{Staging} alias.
    \item Load \texttt{models:/churn\_model/Staging} in a fresh notebook and score a sample, proving the registry round-trip.
\end{enumerate}

\begin{lstlisting}[language=Python,caption={Registering the winning run and loading it back by alias.},label={lst:chapter18-register}]
import mlflow

best_run_id = "<run id from the experiment UI>"
model_uri = f"runs:/{best_run_id}/model"

registered = mlflow.register_model(model_uri, "churn_model")

client = mlflow.tracking.MlflowClient()
client.set_registered_model_alias("churn_model", "Staging", registered.version)

# Round-trip proof: load by alias, never by version
model = mlflow.pyfunc.load_model("models:/churn_model/Staging")
sample = spark.read.table("gold.churn_features").limit(100).toPandas()
print(model.predict(sample.drop(columns=["customer_id", "churned", "as_of_date"]))[:5])
\end{lstlisting}

\subsection{Validation Checklist}
\begin{itemize}
    \item Six runs appear in the experiment with parameters, metrics, and artifacts.
    \item Every run carries \texttt{dataset\_as\_of} and \texttt{code\_commit} tags.
    \item The registered model loads by alias in a fresh notebook and scores.
    \item The experiment UI comparison reproduces your winner selection.
    \item No run references data past its as-of date (Chapter~17's contract).
    \item A teammate can find the best run from the experiment UI without asking you.
\end{itemize}

\section{Friday: Use Case and Architecture}
\index{hyperparameter tuning!collaborative}\index{experiment governance}
A data-science guild runs hundreds of hyperparameter tuning runs per week across five practitioners in one workspace. Without conventions, the experiment list becomes an unsearchable swamp of \texttt{run-473-final-FINAL}. The task: a collaborative tracking workflow.

\subsection{The Workflow}
\begin{itemize}
    \item \textbf{One experiment per problem}, not per person or per day: \texttt{churn-model}, \texttt{demand-forecast}. Runs accumulate comparable history.
    \item \textbf{Mandatory tags}: owner, hypothesis (one sentence: ``deeper trees capture regional effects''), dataset version, code commit. The hypothesis tag is what separates search from wandering.
    \item \textbf{Stable headline metric}: every run logs \texttt{holdout\_auc} (or the program's equivalent), enabling cross-run, cross-week, cross-person sorting.
    \item \textbf{Search discipline}: coarse grid, then narrow; random search over pointless full grids; every configuration tracked, including the failures---negative results are organizational memory too.
    \item \textbf{Registry gates}: promotion to \texttt{Staging} requires the metric threshold and a reviewer; \texttt{Production} requires the Chapter~20 deployment checklist.
\end{itemize}

\begin{table}[H]
\centering
\small
\begin{tabular}{@{}p{4.2cm}p{9.6cm}@{}}
\toprule
\textbf{Element} & \textbf{Convention} \\
\midrule
Experiment naming & \texttt{<problem>-<model-family>}, one per problem \\
Required tags & owner, hypothesis, dataset\_as\_of, code\_commit \\
Headline metric & One stable name per program, logged on every run \\
Promotion gate & Metric threshold + reviewer for Staging; deployment checklist for Production \\
Rollback & Alias move to previous version; rehearsed quarterly \\
\bottomrule
\end{tabular}
\caption{Collaborative tracking conventions.}
\label{tab:chapter18-conventions}
\end{table}

\begin{pitfall}
Comparison across runs is only valid when the data and split are constant. Comparing a run on last month's feature snapshot to one on this month's is measuring data drift, not model quality. The \texttt{dataset\_as\_of} tag is what makes comparisons honest---filter to a common snapshot before declaring a winner.
\end{pitfall}

\section{Operational Guidance for Week 18}
Tracking is cheap; untracked work is expensive. Make the tracked path the easiest path: notebook templates with autolog and tags pre-wired, experiment naming enforced by convention review, and registry promotion as a permission.

\subsection{Performance and Reliability}
Autolog adds negligible overhead for tabular frameworks. Keep artifact sizes sane (log the model, not the training set). Archive stale experiments per retention policy. Rehearse alias rollback: the first production rollback should not be during an incident.

\subsection{Security and Governance Checklist}
\begin{itemize}
    \item Registry write and alias-move permissions restricted to designated roles.
    \item No secrets or personal data in logged parameters, tags, or artifacts.
    \item Dataset tags reference governed table versions, never laptop paths.
    \item Promotion requires the metric evidence attached to the run.
    \item Experiment retention policy documented; organizational memory is curated, not hoarded.
    \item Model cards or equivalent documentation attached at registration.
\end{itemize}

\section{Interview Q\&A}
\begin{enumerate}
    \item \textbf{Q: What does \texttt{mlflow.autolog()} capture per run?}\\
    A: Framework-specific parameters, training and evaluation metrics, and the model artifact itself---automatically---while you add business metrics and metadata tags explicitly.
    \item \textbf{Q: Why register a model instead of referencing run artifacts directly?}\\
    A: The registry makes versions immutable, decouples consumers through aliases, and centralizes the audit trail of promotions---governance that run-artifact paths cannot provide.
    \item \textbf{Q: Which metadata should accompany every experiment run?}\\
    A: Parameters, metrics, the dataset version or as-of date, the code commit, and the environment---the minimum contract for reproducibility and honest comparison.
    \item \textbf{Q: How do you compare hundreds of hyperparameter runs efficiently?}\\
    A: One experiment per problem, one stable headline metric, mandatory tags---then sort and filter in the experiment UI, restricting comparisons to a common dataset snapshot.
    \item \textbf{Q: What does model versioning enable for rollback?}\\
    A: Aliases over immutable versions: rollback is moving \texttt{Production} back to the previous version---instant, reviewed, and reversible, with consumers never changing code.
    \item \textbf{Q: Why must comparisons filter to a common dataset snapshot?}\\
    A: Runs on different snapshots differ in both model and data; without a fixed snapshot you measure data drift and call it model quality.
\end{enumerate}

\section{Further Reading}
Review MLflow in Fabric \cite{microsoftMlflowFabric} and the upstream MLflow documentation \cite{mlflowDocs}. The feature inputs come from Chapter~17's governed feature tables \cite{microsoftFabricDataScience}; the deployment continuation is Chapter~20 \cite{microsoftFabricMlTraining}.

\section*{Key Takeaways}
\begin{itemize}
    \item Experiments group runs; runs record parameters, metrics, artifacts, and tags; the workspace makes them organizational memory.
    \item \texttt{mlflow.autolog()} plus explicit business metrics and metadata tags is the complete tracking contract.
    \item The registry's immutability and aliases are what make promotion and rollback governance instead of archaeology.
    \item Consumers load models by alias; humans move aliases under review.
    \item Collaborative search needs conventions: one experiment per problem, hypothesis tags, a stable headline metric.
    \item Comparisons are valid only on a common dataset snapshot.
    \item Untracked models are debt; tracked, registered, alias-governed models are assets.
\end{itemize}

\section{Exercises}
\begin{enumerate}
    \item \textbf{(Conceptual)} Explain why an immutable version plus a movable alias beats redeployment for rollback.
    \item \textbf{(Conceptual)} What information does a hypothesis tag preserve that metrics alone cannot?
    \item \textbf{(Hands-on)} Extend the grid to 24 runs with random search and compare search efficiency against the coarse grid.
    \item \textbf{(Hands-on)} Add a business metric (expected retained revenue per flagged customer) and re-rank the runs.
    \item \textbf{(Design)} Write the promotion-gate policy: thresholds, reviewers, and evidence required at Staging and Production.
    \item \textbf{(Design)} Two teams share one workspace and keep colliding on experiment names. Design the namespace and the enforcement mechanism.
    \item \textbf{(Interview)} In five sentences, argue to a skeptical engineer that experiment tracking is infrastructure, not bureaucracy.
    \item \textbf{(Operational)} Rehearse a rollback: move \texttt{Production} to the previous version, verify consumers, and document the elapsed time.
\end{enumerate}

\section{Capstone Project: Governed Model Search and Registration}\label{sec:ch18-capstone}
\index{capstone project}\index{MLflow!capstone}\index{model registry!capstone}

\begin{capstonebox}
\textbf{Mission.} Run a governed model search over the Chapter~17 feature table: a tagged, autologged 20-run random search, cross-run comparison on a fixed snapshot, registry promotion through the Staging gate with evidence, and a rehearsed alias rollback---the complete Week~18 discipline in one artifact set.

\textbf{Estimated time.} 4--6 hours.

\textbf{What you'll practice.}
\begin{itemize}
    \item Random search with full tracking and hypothesis tags.
    \item Snapshot-locked comparison and winner selection.
    \item Registry promotion with evidence and review.
    \item Alias-based rollback under a drill.
\end{itemize}
\end{capstonebox}

\subsection*{Scenario}
Contoso's churn program needs its first governed model. The guild's conventions (Friday's table) are adopted; you are the first team to execute them end to end, setting the example the guild will be measured against.

\subsection*{Requirements}
\begin{itemize}
    \item Twenty tracked runs from a random search, each with the four mandatory tags and the headline metric.
    \item A comparison report (screenshot or export) filtered to one \texttt{dataset\_as\_of} snapshot.
    \item \texttt{churn\_model} registered; winner promoted to \texttt{Staging} with the metric evidence attached.
    \item A second version registered and promoted to \texttt{Production}, then rolled back under drill, with timings.
    \item A short model card at registration: purpose, data, metrics, limitations.
\end{itemize}

\subsection*{Implementation Steps}
\begin{enumerate}
    \item Wire the notebook template: autolog, tags, headline metric.
    \item Execute the 20-run search on the fixed snapshot.
    \item Compare, select, and register the winner; attach the model card.
    \item Promote to \texttt{Staging} with evidence; simulate a challenger and promote it to \texttt{Production}.
    \item Execute the rollback drill; document elapsed time and verification.
\end{enumerate}

\subsection*{Deliverables}
\begin{itemize}
    \item The search notebook and the comparison evidence.
    \item Registry evidence: versions, aliases, promotion history.
    \item The model card.
    \item The rollback drill report.
\end{itemize}

\subsection*{Acceptance Criteria}
\begin{itemize}
    \item[\(\square\)] All 20 runs are tagged, metric-bearing, and snapshot-locked.
    \item[\(\square\)] The winner selection is reproducible from the experiment UI by a reviewer.
    \item[\(\square\)] Promotion to each alias carried its evidence and reviewer.
    \item[\(\square\)] The rollback restored the previous version with consumers untouched.
    \item[\(\square\)] The model card states purpose, data, metrics, and limitations.
    \item[\(\square\)] No run logs secrets, personal data, or laptop paths.
\end{itemize}

\subsection*{Stretch Goals}
\begin{itemize}
    \item Add a second model family (gradient boosting) to the same experiment and compare families on the fixed snapshot.
    \item Automate the Staging gate as a notebook asserting the threshold before promotion.
    \item Add per-region metric logging and include a fairness slice review in the gate.
    \item Template the whole flow so the next problem starts at step three, not step one.
\end{itemize}
